\documentclass[12pt]{article}
\usepackage[margin=1in]{geometry}
\usepackage{times}
\usepackage{setspace}
\usepackage{enumitem}
\usepackage{titlesec}
\usepackage{amsmath}
\usepackage{hyperref}
\usepackage{natbib}
\titleformat{\section}{\normalfont\Large\bfseries}{\thesection}{1em}{}
\titleformat{\subsection}{\normalfont\large\bfseries}{\thesubsection}{1em}{}
\titleformat{\subsubsection}{\normalfont\normalsize\bfseries}{\thesubsubsection}{1em}{}

\onehalfspacing

\title{Statement of Research}
\author{Tapas Ranjan Rath}
\date{}

\begin{document}

\maketitle

\section{Planned Independent Research and Publications}

\subsection{Introduction: Human-AI Interaction and the Challenge of Agency}

Technology has become a transformative force in modern society, seamlessly integrating into various aspects of human life and reshaping how individuals live, work, and interact. Among the most significant advancements is artificial intelligence (AI), ranging from smart home assistants to generative AI tools. AI has revolutionized daily routines by enhancing efficiency, convenience, and decision-making processes. Beyond conventional AI, agentic AI represents a new frontier, functioning autonomously to achieve specific goals without continuous human intervention. As these technologies become increasingly embedded in everyday life, they raise critical questions about human-AI interaction, particularly how AI influences individual cognition, decision-making, and the sense of agency.

Traditional metrics for evaluating the success of AI tools—such as sales figures or user engagement—often fail to capture the nuances of user experience, particularly due to an overreliance on behavioral data and the influence of algorithmic biases \citep{NBERw30981, ekstrand2016behaviorism}. Recent research \citep{cornelio2022sense, berberian2019man}, however, suggests that a more comprehensive framework for assessing the impact of AI integration should consider users’ sense of agency (SoA)—the subjective experience of initiating voluntary actions and exerting control over external events.

Sense of agency is a complex, multifaceted phenomenon \citep{pacherie2011self}. Researchers have studied SoA in spatiotemporally proximal and distal perceptual events \citep{kumar2014naturalizing, kumar2017multi}, consistent with the hierarchical event-control approach \citep{jordan2003emergence}. One of the prominent models of SoA is the comparator model \citep{wolpert1995internal, miall1996forward}, which suggests that agency is retrospectively attributed after an action is performed. According to this model, SoA arises from a comparison between sensory feedback and the user’s initial intentions. 
% While SoA has been central in cognitive science, there remains a significant gap in research exploring its role in human-computer interaction, particularly within AI-driven systems.


\citet{cornelio2022sense} categorized human-computer integration systems into three types—body augmentation, action augmentation, and outcome augmentation—each influencing SoA in distinct ways. Body augmentation extends physical capabilities (e.g., prosthetics), while action augmentation supports the user in executing intended actions by compensating for limitations in input command (e.g., touchpads, brain interfaces), motor actuation (e.g., exoskeletons, electrical stimulation), or decision-making via intelligent systems (e.g., autocompletion predictors, autonomous driving, autoplay features). Outcome augmentation, by contrast, modifies the environment to align with user goals (e.g., haptic illusions in VR). Among these, action augmentation—particularly through intelligent systems—remains largely underexplored in relation to SoA.

More recent research proposes a prospective model of SoA \citep{wenke2010subliminal, chambon2012sense, chambon2014action, sidarus2017investigating}, emphasizing the \textit{intention-action} link within the \textit{intention-action-effect} chain. This model suggests that the fluency of action selection itself can serve as a cue to agency, even before an action is executed. This highlights the importance of studying action augmentation through intelligent systems.

\citet{mann1998humanistic} defines intelligent systems as technologies with humanistic intelligence that either \textit{influence user behavior} (e.g., recommenders) or \textit{act on their behalf} (agentic AI). For instance, AI-driven translation tools refine user input while preserving agency. Consider a more complex example: an intelligent alarm system that autonomously adjusts a user’s wake-up time \citep{cornelio2022sense}. Suppose the user has a meeting at 8:00 AM, and the system accesses their calendar, location, and real-time weather data. Detecting overnight snowfall that could cause delays, the system proactively sets the alarm 15 minutes earlier to ensure timely arrival. Here, the AI intervenes in the user’s planned action (when to wake up) while preserving the higher-level goal (being on time), illustrating how intelligent systems can modify actions without undermining agency.

These examples illustrate two distinct forms of action augmentation: action suggestion (e.g., predictive text) and action execution (e.g., automated alarms). Understanding their differential impact on SoA is critical for designing AI systems that respect human autonomy.

\subsection{Research Questions}

Previous researchers have noted the challenge of achieving cognitive coupling between humans and machines \citep{berberian2019man}. \citet{lukoff2021design} found that autoplay and recommendation features can diminish SoA, while others \citep{kumar2014naturalizing, kumar2017multi, ueda2021influence} observed enhanced SoA with automation. However, most prior work has either focused on the action-effect link in the chain or examined SoA in categories and sub-categories other than intelligent systems, as \citet{cornelio2022sense} outlined.


I aim to systematically study the impact of both types of intelligent systems—action recommenders and action executers—on users’ sense of agency (SoA), with a particular focus on the \textit{intention–action} link within the \textit{intention–action–effect} chain.
Specifically, the proposed questions are:

\begin{enumerate}[label=\arabic*.]
    \item How do AI-driven recommendations affect a user’s sense of agency?
    \item How does agentic AI that autonomously executes actions influence SoA compared to user-initiated actions?
    \item How does the user’s sense of agency evolve as their avatar in a virtual game environment learns and adapts over time?
    \item How does user SoA evolve when interacting with intelligent avatars in multiplayer online games?
\end{enumerate}

\subsection{Methodology}

To address these questions, a series of experimental paradigms will be employed that manipulate action selection fluency while measuring both explicit (e.g., self-reports) and implicit (e.g., timing, behavior) indicators of SoA. The methodologies include:

\begin{itemize}
    \item Behavioral paradigms involving multi-alternative choice with AI suggestions varying in automation and personalization.
    \item Controlled virtual game environments in multi-user settings to observe cooperation vs. competition dynamics and their effects on SoA
    \item Computational modeling to predict SoA outcomes in AI-assisted environments.
\end{itemize}

\subsection{Publication Goals}

I aim to publish my findings in leading journals across cognitive science, AI ethics, and human-computer interaction, such as Cognitive Science, AI and Ethics, and International Journal of Human-Computer Studies. Additionally, I plan to contribute to premier conferences such as the ACM CHI Conference on Human Factors in Computing Systems, the Cognitive Science Society Annual Meeting (CogSci), and the ACM Conference on Intelligent User Interfaces (IUI). These venues will ensure that my research reaches interdisciplinary audiences and fosters meaningful discussions on human-AI integration and the sense of agency.

\section{Resources Needed}

To carry out this research, the following resources may be required:

\begin{itemize}
    \item Collaboration with AI experts for developing intelligent decision-making systems.
    \item Collaboration with blockchain experts for secure, decentralized game environments.
    \item Access to human-computer interaction laboratories for experimental work.
    \item Computational resources for simulations and machine learning model development.
\end{itemize}

\section{Building on and Extending Previous Work}

My current research explores context effects in decision-making, particularly how adding a third inferior option (attraction effect) influences choice. This work has highlighted how choice difficulty contributes to decision biases.

The proposed research extends this work in several indirect ways:

\begin{itemize}
    \item From context-based biases to AI-driven decision biases and agency: Investigating how AI recommendations and autonomous actions influence SoA and introduce new biases.
    \item From choice difficulty to action selection fluency: Studying how fluency in selecting an action affects SoA in AI-assisted contexts.
    \item From static decision-making to dynamic AI interaction: Exploring how AI learning dynamically alters user agency in iterative environments.
\end{itemize}

This research will systematically investigate how AI-driven recommendations and autonomous actions influence SoA, offering new insights into AI ethics, human-computer interaction, and cognitive psychology. By integrating cognitive theories of agency with human-computer interaction research, I aim to design AI systems that enhance user experience while preserving autonomy and agency.

\section{Conclusion}

As AI continues to influence human decision-making, understanding its effects on agency is critical. This research aims to advance theoretical models of agency and inform the design of AI systems that support autonomy, ethical interaction, and user empowerment.

\bibliographystyle{apalike}
\bibliography{references} % Ensure you have a references.bib file
\end{document}